\section{Introducción}
\label{sec:introduccion}

\subsection{Continuidad del proyecto}

El Sistema de Gestión de Solicitudes de Soporte Técnico de Internet es el proyecto final de
carrera del equipo ACC para la asignatura de Aplicaciones Distribuidas. En la Entrega 1 se
definió el dominio del problema —un proveedor de servicios de Internet (ISP) que necesita
registrar, clasificar, asignar y dar seguimiento a solicitudes de soporte técnico de sus
clientes— y se construyó un primer prototipo monolítico con comunicación por \emph{sockets}. En
la Entrega 2 el sistema se descompuso en cinco microservicios independientes
(\texttt{auth-service}, \texttt{ticket-service}, \texttt{ai-service}, \texttt{notification-service}
y \texttt{report-service}), comunicados mediante REST y eventos asíncronos sobre Apache Kafka, con
autenticación y autorización por rol basada en JWT~\cite{jones2015jwt} siguiendo los principios de
diseño orientado a recursos de Fielding~\cite{fielding2000rest} y las prácticas de descomposición
de servicios descritas por Newman~\cite{newman2021microservices}.

La Entrega 3, cubierta por este documento, profundiza en dos dimensiones que las entregas
anteriores dejaron abiertas: (i) el comportamiento de la capa de datos bajo fragmentación
horizontal real y su tolerancia a fallos de nodo, y (ii) el procesamiento analítico paralelo de
volúmenes de datos que superan lo que un solo proceso puede manejar con margen razonable de
tiempo. Ambas dimensiones se abordan sobre la misma base de código de la Entrega 2, no como un
sistema nuevo: los cambios de esquema, el rediseño de \texttt{ticket-service} y el pipeline de
Spark se integran con los microservicios y el flujo de eventos ya existentes.

\subsection{Objetivos de la Entrega 3}

\begin{itemize}[leftmargin=1.5em]
    \item Fragmentar horizontalmente la tabla de mayor cardinalidad del sistema
        (\texttt{tickets}) sobre un clúster real de CockroachDB de tres nodos, aplicando el
        criterio de fragmentación exigido para el equipo ACC (\texttt{fecha\_apertura}), y
        verificar formalmente su corrección.
    \item Medir y documentar el comportamiento del clúster ante la pérdida de uno y dos nodos,
        incluyendo latencia de escritura/lectura y disponibilidad efectiva del servicio.
    \item Instrumentar el servicio de datos con métricas operativas (Prometheus/Micrometer) y
        pruebas de integración reproducibles contra una instancia real de la base de datos
        (Testcontainers), evitando que la corrección del sistema dependa únicamente de
        inspección manual.
    \item Diseñar e implementar un pipeline de procesamiento paralelo sobre Apache Spark que
        aplique las cinco categorías de transformación exigidas por la guía de la asignatura
        sobre un conjunto de datos de al menos 500{,}000 registros, orientado a un objetivo
        analítico propio del dominio del equipo: reincidencia de clientes y agrupamiento de
        incidencias.
    \item Diseñar y ejecutar un protocolo experimental formal que contraste el rendimiento del
        pipeline paralelo contra un baseline secuencial equivalente, con análisis estadístico de
        significancia, y que caracterice la ley de Amdahl~\cite{amdahl1967validity} observada en
        el propio pipeline.
\end{itemize}

\subsection{Estructura del documento}

La Sección~\ref{sec:trabajos-relacionados} sitúa las decisiones de diseño de este trabajo frente a
la literatura de bases de datos distribuidas y procesamiento paralelo de datos. La
Sección~\ref{sec:fundamento-teorico} resume los fundamentos teóricos aplicados: teoremas de
fragmentación, consistencia distribuida y las leyes de escalabilidad paralela. La
Sección~\ref{sec:diseno-esquema} documenta el rediseño del esquema y la política de fragmentación.
La Sección~\ref{sec:cluster-verificacion} presenta la verificación empírica del clúster: pruebas
de integración, métricas e instrumentación, y los resultados de tolerancia a fallos. La
Sección~\ref{sec:pipeline-paralelo} describe el pipeline de Spark y sus cinco transformaciones. La
Sección~\ref{sec:protocolo-resultados} presenta el protocolo experimental y sus resultados. La
Sección~\ref{sec:discusion} discute las limitaciones y \emph{trade-offs} del diseño, y la
Sección~\ref{sec:conclusiones} cierra con las conclusiones y el trabajo futuro.
